\section{Brainwashing and a Liberal Education}

Yesterday afternoon, while I was finishing up a bottle of wine and waiting for the lasagna to bake, I did some channel surfing. I found a channel called EWTN with an upcoming news show from a Catholic perspective; that sounded better than the secular news shows. Although I never watch that channel, I am familiar with it. I know, for example, that a popular show was canceled when the priest was discovered to be living a secret life of pleasure and luxury. Another show was canceled when the host blamed boys for seducing priests when they are vulnerable. Normally, I follow Terence's dictum that ``nothing human is alien to me", but that comment is alien to me.

The news show began with a discussion of the political situation in Egypt; an Egyptian Christian was the guest. Of course, the first topic was the situation of Christians in Egypt, which apparently is precarious. It would have been a good opportunity to discuss how the Muslims and Copts had got along for the past 1500 years, but that was not mentioned. However, the show took an interesting turn, at least for me, although the two interlocutors were oblivious. They mentioned a few Muslim radicals who had undergone an education in the West. They were genuinely surprised how a liberal Western education did not ``take" in these men, who nevertheless pursued quite illiberal policies.

While, in their eyes, that was a perfectly natural attitude to take, for me, it was quite jarring. Why bother to watch news from a ``Catholic" perspective if it is indistinguishable from secular sources? Perhaps those Muslims, as outsiders, were immune to the implicit supposition that the dogmas of liberalism are rational and axiomatic. \textbf{Mohamed Morsi} got a graduate degree in materials science at USC, a STEM major, so he would have missed out on most of the liberal indoctrination. Apparently, even California dreaming did not suffice to turn him.

So what is the attraction of liberal democracies to the religious news hosts? I have mentioned Pope Pius IX's \textit{Syllabus of Errors}\footnote{\url{http://www.ewtn.com/library/papaldoc/p9syll.htm}} several times and now is the time to look at its points. It has earned the approval of \textbf{Julius Evola} for its adamant opposition to liberalism. Now it is true that the current pope claims the syllabus is now overturned, as there is a different critique of liberalism. If so, then what is it? Not even the news anchor is aware of it. Keep in mind that Evola never rejected the \textbf{Ancien Regime} even if his views on sexual matters were looser. In reading these propositions, readers may, if they prefer, substitute their own spiritual authority for any mention of the Church.

The literary form of the Syllabus is to state a positive proposition while asserting it as false. Point 3 is the fundamental dogma of Enlightenment liberalism:

\begin{quotex}
3. Human reason, without any reference whatsoever to God, is the sole arbiter of truth and falsehood, and of good and evil; it is law to itself, and suffices, by its natural force, to secure the welfare of men and of nations. 

\end{quotex}
We hope that EWTN agrees about the falsity of (3). This aligns them more with Morsi than with liberal education. Point 39 is related:

\begin{quotex}
39. The State, as being the origin and source of all rights, is endowed with a certain right not circumscribed by any limits. 

\end{quotex}
Clearly the Western Tradition is in agreement with Morsi on this point. The next is more controversial:

\begin{quotex}
40. The teaching of the Catholic Church is hostile to the well-being and interests of society. 

\end{quotex}
In the case of Egypt, readers can substitute ``Islam" for ``Catholic Church". Leaving aside for the moment the historical conflict between Traditions in geographic proximity, liberalism rejects any and all spiritual authority. That is why any attempts to ``modernize" or to be acceptable to the modern world are ultimately futile. The next proposition is asserted by the French and American Revolutions. With the appropriate change of language, this proposition is also rejected by Morsi.

\begin{quotex}
55. The Church ought to be separated from the State, and the State from the Church. 

\end{quotex}
Proposition 77 is basic for the liberal state as well as by Vatican II under the rubric of ``freedom of religion". Neo-pagans also accept prop 77, thus aligning themselves with liberalism rather than Tradition. We have pointed out that the Ancient City did not accept religious freedom except in an attenuated form. With a suitable verbal transposition, the Islamic state also rejects proposition 77:

\begin{quotex}
77. In the present day it is no longer expedient that the Catholic religion should be held as the only religion of the State, to the exclusion of all other forms of worship. 

\end{quotex}
This is just a sample. The point is that we are hard pressed to determine how the contemporary so-called Catholic viewpoint differs from what is absorbed in a ``liberal education". If we turn to EWTN, we certainly don't learn much. We can end with an issue of current relevance. One of the functions of the spiritual authority is to determine what is ``real", a prime example, that of the ``Real Presence". Another is what constitutes a ``real", rather than merely a nominal, marriage. We see that the same issue was debated 150 years ago, even if the specifics differ. This false proposition is:

\begin{quotex}
68. The Church has not the power of establishing diriment impediments of marriage, but such a power belongs to the civil authority by which existing impediments are to be removed. 

\end{quotex}
In other words, the spiritual authority alone determines the impediments to marriage. If left to the temporal powers alone, the result will always be wrong. Hence, attempts to ``debate" in our time what constitutes marriage are doomed to fail ultimately; the modern mind is not bound by the real but by the nominal.



\flrightit{Posted on 2012-12-03 by Cologero }

\begin{center}* * *\end{center}

\begin{footnotesize}\begin{sffamily}



\texttt{Jason-Adam on 2012-12-03 at 23:17 said: }

This was a very good article, Cologero. 

Sadly I think we all know that the sect based on the Vatican II teachings calling itself the Catholic Church is the direct opposite of Catholicism.

I am sure you've read Fr Rama Coomaraswamy's book ? 

PS – the experience of the shallow hypocrisy that is liberal

education is what led me to reject modernity as a whole


\hfill

\texttt{Michael on 2012-12-04 at 19:39 said: }

My guess is that the EWTN host would have agreed with everything in the Syllabus, but one can't just replace every instance of ``Catholic" with ``Islam" and have the statements still be accurate. Even if the Church insists that the ``Catholic religion should be held as the only religion of the State," there is no one who would listen in the post-modern age. The Church has no choice but to try to work with the liberal culture. The good news (or the bad news) is the attempt to engage the liberal culture is failing miserably. This reminds me of the following quote from Belloc:

``But if I be asked what sign we may look for to show that the advance of the Faith is at hand I would answer by a word the modern world has forgotten: Persecution. When that shall once more be at work it will be morning."

It seems morning is coming soon.

I know that this blog is about Tradition, not theology, but logic still holds. Either Jesus was God or he was not. If he is God, then he built his Church that we must obey. We have his promise that the gates of hell will not overcome his Church so we don't have the option of leaving it simply because large portions of it have capitulated to the liberal culture. There are plenty of Catholics who remain faithful to the Church as she has stood for 2000 years. ``Yet I have left me seven thousand in Israel, all the knees which have not bowed unto Baal."

As for Islam, I understand that there are certain aspects of the culture that reflect Tradition, but it is only by derivation. Would you really want to belong to the culture that Morsi is implementing? And even if you prefer it to the Church, is it an option if the Jesus is God?


\hfill

\texttt{Jason-Adam on 2012-12-04 at 22:37 said: }

I would answer that Islam is still preferable to a godless society.

The Church does not have to engage the world, if the world is refusing to follow God then the Church must withdraw from the world and return to the catacombs.

Indeed, despite all the attempts of Satan, God's Church is still here and is destined to prevail in the End…….


\hfill

\texttt{Cologero on 2012-12-05 at 00:06 said: }

Given recent comment topics, here is a link to Roger Garaudy's conversions from Marxism to Christianity to Islam: \textit{The philosophical itinerary of Roger Garaudy}\footnote{\url{http://rogergaraudy.blogspot.co.uk/2011/02/philosophical-itinerary-of-roger.html}}. It may be of interest, or maybe not.

I don't know what the Sheik would have thought of Ibn Arabi, Hafiz, or Omar Khayyam who all wrote poems of love. Perhaps Shia is more in tune with the mentality of the West.


\hfill

\texttt{Jason-Adam on 2012-12-05 at 17:18 said: }

I would say, from reading the article, that Garaudy's adoption of Islam was due to residual anti-Western-ness from his Communist past. Many people, even after they realise their need for transcendence in their lives, have still too much ``liberal guilt" and self-hatred towards their own people and culture to look for God in the Western past.

I would say also that the elements of Shia that are akin to the Western mind are simply Islamic borrowings from Catholicism.


\hfill

\texttt{jAh reZident on 2012-12-06 at 17:56 said: }

Surely Tradition points to the way of the One God, and emanates from a received wisdom, ultimately delivered by the prophets Moses, Jesus and Mohammed (pbuh). So `adoption' of Islam is a misnomer, Jesus as a `god' is a reversal of tradition, and `modernity' is the ploy of the Enemy to prepare the way for the dominance of all that opposes the Real.


\hfill

\texttt{Dominion on 2012-12-07 at 21:51 said: }

Would you elaborate? I gather you might be thinking of things like devotion to Sufi saints and a greater focus on things like devotional practices and pilgrimages, but I'd say these things arose spontaneously during the transition to Islam, built on local practices, despite their similarity to Catholic practices. Both Shi'ite Iran and Catholic Europe built upon Indo-European religious bases, so it's not too unexpected.


\hfill

\texttt{Alexander Maximilian on 2012-12-08 at 14:21 said: }

There is an excellent book called ``EWTN a Network Gone Wrong" about how the have become a station trying to compete with other stations and thereby have sold out to modernism a lot if it due to the number of Protestant converts that produce many of the shows.


\end{sffamily}\end{footnotesize}
